\section{Preregistered Evaluation}
\label{sec:experiments}

The study is frozen before confirmatory search and answers three primary
questions: Does mechanism-visible search outperform a score-only optimizer
under an otherwise identical budget?  Does it improve over the unmodified
harness?  Does it increase verified repair on controlled faults?  The
corresponding minimum effects are preregistered rather than selected after
unsealing.

\subsection{Conditions and matched budgets}

\begin{table}[t]
  \caption{Experimental conditions.  \scorearm{} is the causal control for
  mechanism visibility; prompt-only is a scope baseline, not a substitute for
  that control.}
  \label{tab:conditions}
  \centering
  \begin{tabular}{p{1.05in}p{4.00in}}
    \toprule
    \textbf{Condition} & \textbf{Authorized search behavior} \\
    \midrule
    \staticarm & Frozen seed harness; no search. \\
    Random-valid & Uniform finite-catalog mutation under the same stopping and
      execution ceilings. \\
    Prompt-only & Optimizer edits prompts only; its feedback view is frozen and
      disclosed. \\
    \scorearm & Full mutation grammar; aggregate score, uncertainty, resource
      totals, and redacted gate decision only. \\
    \fullarm & Same grammar and budget as \scorearm, plus source-linked
      activation, adherence, behavior, revert, and placebo evidence. \\
    \bottomrule
  \end{tabular}
\end{table}

Within every model--task--seed block, \scorearm{} and \fullarm{} share the seed
harness, proposer and sampling settings, editable surfaces, proposals per
round, solver, task and rollout coordinates, nomination rule, number of gate
queries, and hard resource ceilings.  Cache state is neutralized or charged
symmetrically.  All proposer, diagnoser, solver, grader, retry, and failed-call
usage is retained.

\paragraph{External harness optimizers.}
Secondary comparisons include a history-rich Meta-Harness-style optimizer
\citep{lee2026metaharness} and a same-model weakness-mining/regression-gate
optimizer based on the published Self-Harness protocol
\citep{zhang2026selfharness}, provided their executable releases and licenses
permit the selected task.  An official release is named as such; missing final
task code is replaced only by an independently implemented
\emph{paper-algorithm reconstruction} and is never described as an official
reproduction.  These arms use the same solver, task coordinates, and disclosed
hard budget, but remain secondary because their mutation and feedback spaces
cannot in general be made identical to the primary \fullarm/\scorearm{}
contrast.

\subsection{Models, tasks, and independent runs}

The frozen roster uses
\designvalue{three open model configurations spanning at least two families}
and \designvalue{four task families}: mathematical generation, multi-type
reasoning, high-label classification with group separation, and strict
structured generation or coding with source/time separation.  Availability
smoke tests may select deployable endpoints but may not inspect benchmark
scores.  Every arm within a cell uses the same exact model revision and
inference configuration.

One complete search is the repeated algorithmic unit.  The number of independent
search seeds is \evidenceblocked{pilot variance estimates and the frozen power
simulation are incomplete}.  The simulation selects from a predeclared roster
and runs the exact multiplicity-aware analysis; rollout repeats do not count as
independent search replications.

\subsection{Estimands and analysis}

For fixed model roster \(\mathcal{M}\) and task roster \(\mathcal{T}\), the
primary real-task estimands weight each model--task cell equally:
\begin{align}
\Delta_{\mathrm{causal}}
  &= \frac{1}{|\mathcal{M}||\mathcal{T}|}
     \sum_{m,t}\left[\mu(\text{\fullarm},m,t)
       -\mu(\text{\scorearm},m,t)\right],\\
\Delta_{\mathrm{e2e}}
  &= \frac{1}{|\mathcal{M}||\mathcal{T}|}
     \sum_{m,t}\left[\mu(\text{\fullarm},m,t)
       -\mu(\text{\staticarm},m,t)\right].
\end{align}
The third estimand is the \fullarm{}-minus-\scorearm{} difference in verified-repair
probability on HarnessFaultBench.  The preregistered minimum effects are 0.02
absolute normalized score for each real-task contrast and 0.15 for verified
repair.  Three one-sided tests use Holm--Bonferroni family-wise
\(\alpha=0.05\) \citep{holm1979simple}.  A headline conclusion requires all
three adjusted tests to pass and all point estimates to exceed their minimum
effects.

Search seeds are resampled within model--task cells and task groups are
resampled within runs.  A robust hierarchical paired model is primary; an
equal-cell studentized cluster bootstrap is the frozen fallback if it fails to
converge.  Failure-as-zero is the intention-to-treat analysis, accompanied by
complete-case and worst-case sensitivity analyses.  Individual model, task,
surface, and factor findings are corrected or labeled exploratory.

\subsection{Selection closure, sealing, and cost}

Search stops at the first frozen round, proposal, nomination, optimizer-call,
resource, patience, or invalid-candidate limit.  All configured cells,
including failures, must close before the sealed authority accepts champion
hashes.  It evaluates the entire grid in one release.  A post-unseal change
starts a new preregistration and new sealed lineage.

We report score and protected outcomes alongside uncached and cached input,
output/reasoning tokens, normalized frozen-price cost, actual billed cost when
available, wall time, tool calls, failures, and retries.  Score per million
tokens and a Pareto frontier complement absolute outcomes; cost is not hidden
inside the optimized objective.
